Coordinated cryptocurrency pumps provide a unique environment to study how public messages influence market prices within very short timeframes. In these schemes, a pump signal is broadcast to a large number of traders, instructing them to buy a specific asset simultaneously. If enough participants react quickly, the sudden wave of demand can drive the price up for a brief period, allowing early buyers or insiders to sell at a profit to those who enter later. The core question for this study is to identify which observable elements of these signals help explain the short-term price reaction and whether the immediate post-reveal outcome can be forecast from information available before the reveal.

This setting is particularly suitable for research because the signals are public, the timing is precise, and the market's response follows almost immediately after the reveal. Messaging platform pump groups typically use countdowns, forceful wording, and explicit trading instructions to organize their members. These messages often lack any new fundamental information about the asset. Instead, their primary function is to coordinate attention and trading pressure at a single moment. Consequently, this thesis treats coordinated signals as tools for social coordination rather than as conventional financial news.

This chapter reviews the literature that explains why such signals can matter. It begins with market manipulation and cryptocurrency pump-and-dump schemes, then moves to communication, social coordination, market microstructure, and short-horizon price reactions. It then separates detection and prediction studies from causal measurement. Finally, it develops the hypotheses and explains why the empirical analysis must keep causal estimation and prediction as two related but separate tasks.

\section{Cryptocurrency Market Manipulation and Pump-and-Dump Schemes}
\label{sec:crypto-manipulation-literature}

Market manipulation has been a persistent challenge in financial history. Theoretical models suggest that prices can be moved when traders take actions that others interpret as valuable information \cite{allen1992stock}. Manipulators do not always require a deep understanding of an asset's intrinsic value. Instead, it often suffices to generate a surge in buying pressure or to lead other participants to believe that demand is rising. A pump-and-dump scheme follows this logic, relying on the combined effect of price pressure and informational signals to move the market.

Empirical studies of traditional stock markets show that manipulation is most common where there are significant gaps in information access, timing, or control. Manipulation often involves trading activity that creates artificial price movements, followed by the manipulator's profitable exit \cite{aggarwal2006stock}. Research on emerging markets also highlights how weak oversight and the presence of powerful intermediaries can make price manipulation easier to carry out \cite{khwaja2005unchecked}. These findings are highly relevant to cryptocurrency markets, where organized pump groups have been observed across many coins, channels, and exchanges \cite{kamps2018pump, hamrick2021examination}.

The structure of cryptocurrency markets introduces additional factors that facilitate coordinated manipulation. Trading is fragmented across numerous exchanges, leading to price discrepancies across different markets \cite{makarov2020trading}. Furthermore, many digital assets suffer from extreme volatility and inconsistent trading volume \cite{liu2021risks}. Evidence from centralized exchanges also shows that wash trading can inflate reported volume and temporarily distort prices on unregulated venues \cite{cong2023crypto}. Some tokens trade in such "thin" markets that even a modest wave of buying can move their price much more than it would in a more liquid market. While these features do not mean that every sharp price change is the result of manipulation, they do suggest that coordinated buying can have a disproportionate impact when market depth is low.

The early history of cryptocurrency trading also shows that manipulation is not a marginal issue. Research showed that suspicious trading patterns played a significant role in Bitcoin's early price history \cite{gandal2018price}. Related evidence from the 2017 boom comes from \cite{griffin2020untethered}. The study examines whether Tether flows helped support Bitcoin prices and reports that purchases made with Tether often occurred after Bitcoin market downturns and were followed by large price increases. This pattern is presented as evidence consistent with price support, not as proof that ordinary demand played no role. These findings are important because they show that manipulation can appear even in the best-known cryptocurrency market, not only in small and obscure tokens. The broader lesson is that cryptocurrency prices can be shaped by trading behavior, market design, and information gaps, not only by long-term beliefs about technology.

Organized pump-and-dump schemes became one of the clearest forms of this problem. A pump-and-dump scheme normally has two connected stages. In the pump stage, organizers create attention and buying pressure around a selected asset. In the dump stage, early buyers or insiders sell into the demand created by later participants. The price increase is therefore artificial and short-lived. It is not based on a new fundamental valuation of the project, but on coordinated timing.

More recent research has focused specifically on organized pump groups. Methods for defining and detecting these events were introduced by \cite{kamps2018pump}, while \cite{victor2019cryptocurrency} later measured confirmed Telegram/Binance pump events and tested a detection model. A broad ecosystem study showed that Telegram and Discord groups coordinated thousands of pumps targeting hundreds of coins, and that less popular coins experienced stronger price increases \cite{hamrick2021examination}. Other studies have mapped the internal structure of messaging groups and linked public signals to specific market reactions \cite{xu2019anatomy, li2021cryptocurrency}.

This body of research can be organized into several related branches. The first branch defines the event itself and develops rules for detecting suspicious price and volume movements. This work is necessary because pump-and-dump schemes are not always announced in a clean or standardized form, and market data alone can contain many sharp movements that are not manipulative \cite{kamps2018pump, victor2019cryptocurrency}. A second area studies the anatomy of online pump groups. It shows how organizers prepare their members, reveal the target, and try to manage the early buying wave after the announcement \cite{xu2019anatomy, lamorgia2023doge}. A third branch takes an ecosystem view and shows that these events are not rare accidents, but repeated activities across many assets, channels, and exchanges \cite{hamrick2021examination, li2021cryptocurrency}.

These areas are closely connected, but they answer different questions. Detection studies ask whether a suspicious event can be found. Anatomy studies ask how a group organizes the event. Ecosystem studies ask how large and widespread the problem is. Target-prediction studies ask whether the selected coin can be guessed before the public reveal \cite{hu2023sequence}. Real-time detection studies then ask whether suspicious behavior can be identified quickly enough to warn market participants \cite{lamorgia2020pump, lamorgia2023doge}. Together, these studies show that pump-and-dump schemes leave observable traces in both messages and market data.

The focus of this thesis is narrower. It does not mainly ask whether pump-and-dump schemes exist or whether they can be detected. It asks whether differences in the public signal help explain the immediate market response once the target coin is revealed. This requires moving from event identification to signal interpretation. A pump announcement is not only a label marking the start of an event. It is also the moment when attention, timing, and practical trading instructions are compressed into one public message. For that reason, the wording and structure of the signal can be studied as part of the manipulation mechanism.

This distinction also helps explain why communication variables should not be treated as minor details. A group can select the same coin but present the signal in different ways. One message may be short and urgent, another may include longer instructions, while another may follow a stronger buildup of preparatory hype. These differences may shape how quickly members understand the message and whether they believe other members are about to act. Existing pump-and-dump research provides the foundation for this question, but the thesis extends the literature by focusing on variation between signals after the event has already been identified.

These studies on pump-group detection, group organization, and ecosystem scale make cryptocurrency pump-and-dump schemes different from many older forms of manipulation. The manipulation is not hidden behind a complex corporate story or a long public campaign. In many cases, the manipulative purpose is openly visible to group members \cite{dhawan2023new}. The message tells participants to buy quickly, and the group understands that the price move depends on collective action \cite{xu2019anatomy}. This makes the cryptocurrency setting especially useful for studying manipulation as a coordination problem.

The scale of the problem is also important. Pump-and-dump activity has been observed across many channels, exchanges, and coins, not only in a few isolated cases \cite{hamrick2021examination, li2021cryptocurrency}. This repeated structure allows researchers to move beyond case studies. It becomes possible to compare many events and ask why some signals produce stronger market reactions than others. That comparison is central for understanding whether communication features, such as urgency or instruction clarity, matter after controlling for the market conditions around the target asset.

Successful pump-and-dump schemes rely heavily on timing \cite{kamps2018pump, xu2019anatomy}. Typically, organizers select a target asset that is prominent enough to attract interest but small enough that its price can be easily moved \cite{hamrick2021examination, charfeddine2024drives}. Prior to the actual reveal, the group builds anticipation through reminders and countdowns. When the target coin is finally announced, members are encouraged to buy immediately. If the response is fast enough, the price spikes within the first few minutes. As early buyers start to sell, those who entered later are often left holding assets that rapidly lose their value \cite{li2021cryptocurrency, dhawan2023new}.

The literature also suggests that participating in a pump is not always the result of a simple misunderstanding. Many participants are aware of the risks but gamble on being able to exit before the crash. Research on cryptocurrency manipulation describes this as a modern evolution of manipulation, where open participation continues even when the average trader faces poor odds \cite{dhawan2023new}. This perspective is important because it moves the focus away from "fake news" and toward the mechanics of fast coordination in a high-stakes environment.

Taken together, these studies show that cryptocurrency pump-and-dump schemes combine three conditions that rarely appear so clearly at the same time. The communication is public within the group, the timing is narrow, and the target market is often thin enough to react sharply. This combination makes the setting useful for an empirical study, since the message, the event time, and the market response can be linked within one observation.

The role of insiders is also central. Pump organizers know the target asset before ordinary group members do. This creates a timing advantage. Insiders can accumulate the coin before the public reveal, while later participants only receive the signal once the buying race has already started. Evidence of pre-announcement price run-ups and post-announcement reversals supports the view that pump-and-dump schemes can transfer wealth from outsiders to better-positioned insiders \cite{li2021cryptocurrency}. This is one reason why the public signal itself is not a neutral announcement. It is part of a market event in which information, timing, and trading position are unevenly distributed.

Market size also plays a fundamental role. Larger assets can usually absorb significant buying volume without sharp price swings, while smaller assets with limited trading depth are more sensitive to coordinated demand. This relationship is closely tied to illiquidity, since trades have a larger price impact in thin markets \cite{amihud2002illiquidity}. Evidence from pump-and-dump data also shows that less popular coins can react more strongly to these events \cite{hamrick2021examination}.

This logic connects the existing literature to the empirical design of this study. Organizers need coins that are visible enough to attract participants but still "moveable" enough for coordinated demand to affect prices. Market capitalization is therefore useful later as a practical proxy for this vulnerability.

Effective pump-and-dump schemes are also built on communication. In traditional markets, rumors and tips were typically spread through personal networks, newsletters, or early online message boards. Rumors can drive market behavior when traders act before information is fully verified \cite{vanbommel2003rumors}. Research on internet stock boards also shows that social discussion can influence trading even without fundamental news \cite{sabherwal2011internet}. These studies help explain why public signals can matter even when they do not reflect an asset's real value.

Digital messaging platforms have changed the speed and structure of this process. Telegram channels and Discord groups can reach thousands of users at once \cite{hamrick2021examination, lamorgia2023doge}. They also make timing easier to control through scheduled messages, pinned posts, countdowns, and repeated reminders. These tools create a shared moment of attention before the target coin is revealed, so the signal tells participants not only what to buy, but also when and how quickly to act.

A coordinated pump signal is distinct from typical investment discussion. It is usually brief, direct, and focused on immediate action. Rather than encouraging a long-term view of an asset, it demands a rapid response to capture short-term gains before others do. This emphasis on speed makes the phrasing of the signal critical. For example, a message that conveys urgency and clarity is likely to trigger a faster response than one that is vague or poorly timed.

Research on social media and messaging platforms supports this view. Online signals are often connected to suspicious market volatility \cite{mirtaheri2021identifying, xu2019anatomy}, and Telegram or Discord groups are important venues for coordinated pump activity \cite{hamrick2021examination, lamorgia2023doge}. The content and delivery of a message are therefore part of the event, not only background details.

Pump communication usually develops in stages. Prior to the reveal, messages may build anticipation by announcing the event time, warning members to prepare exchange accounts, or using language that suggests a rare opportunity. These messages can increase attention even before the target coin is known. At the reveal, the communication becomes more direct. The target asset, exchange, trading pair, and buying instruction may be placed in one short message. After the reveal, follow-up messages may try to keep buying pressure alive or delay selling by telling members to hold \cite{xu2019anatomy, lamorgia2023doge}.

This staged communication is one reason why the signal cannot be reduced to the coin symbol alone. The target symbol tells participants what to buy, while the surrounding messages shape how the event is understood. A countdown can make the event feel scheduled and shared, an urgent call can create pressure to act immediately, and clear exchange or pair information can reduce practical confusion. These details matter because pump participants operate under high time pressure.

Messaging platforms also create a visible hierarchy between organizers and ordinary participants. Organizers control the channel, choose the reveal time, and decide how much information is released before the event. Ordinary members receive the target only when the public signal appears. This difference in timing is central to the logic of the scheme, since those with earlier knowledge can prepare or accumulate before the rest of the group reacts \cite{xu2019anatomy, li2021cryptocurrency}. The public message is therefore both a coordination device and a point where unequal information becomes visible.

This hierarchy also means that a Telegram channel is more than a technical source of messages. It can represent a repeated organizer style. Some groups may build stronger anticipation, reveal targets more directly, choose narrower assets, or use more aggressive language than others. The same channel may also develop a reputation among its members, even if that reputation is based on past hype rather than reliable performance. Prior research shows that pump activity is concentrated across organized groups and that communication patterns are part of how these events are structured \cite{hamrick2021examination, xu2019anatomy, lamorgia2023doge}. For this reason, channel-level differences matter for interpretation. They may capture organizer habits, audience size, target-selection routines, or repeated message styles that are difficult to observe directly.

This point is important for both causal and predictive analysis. If one channel regularly chooses more moveable coins and also writes more urgent messages, a simple comparison could confuse the channel's target-selection strategy with the effect of urgency itself. At the same time, the same channel pattern may be useful for prediction, since repeated organizer behavior can leave stable traces in the data. The literature therefore supports treating organizer and channel context as part of the event environment, while still keeping the public signal as the main object of interest.

The pivotal moment in any pump is the public reveal of the target coin, denoted as $T_0$. Conceptually, messages around this point matter through three broad coordination dimensions, including how attention is concentrated prior to the reveal, how urgent the reveal appears, and how clearly the signal translates attention into action.

The first dimension is attention concentration. A public signal can matter even when it does not add fundamental information if many traders treat it as a common point of reference. Research on public information shows that public signals can influence behavior when people care not only about their own information, but also about how others are likely to act \cite{morris2002social}. This is directly relevant to pump groups. A member may buy not because the coin has become more valuable, but because the member expects many other members to buy immediately after the reveal. Attention therefore becomes valuable because it creates a shared moment in which traders watch the same asset at the same time.

Attention is also important in retail trading more broadly. Investors often buy assets that first catch their attention, especially when they face too many possible assets to search through carefully \cite{barber2008all}. Pump groups exploit this attention problem in a very direct way. They remove the search process by naming one coin and presenting it as the focus of the group. Preparatory messages can make this effect stronger by increasing anticipation before the reveal. In the empirical design of this thesis, \textit{Hype Intensity} captures this attention-building process by measuring how strongly the preparatory communication concentrates near $T_0$.

The second dimension is urgency. Financial text can matter because tone and wording shape how traders interpret market situations. Evidence from media-based finance research shows that textual content can be linked to returns and trading volume, even when the text is not simply a direct report of new fundamental information \cite{tetlock2007giving}. In pump groups, urgency is more immediate than ordinary market sentiment. Urgent wording tells members that delay is costly. It narrows the time available for reflection and makes the event feel like a race against other participants. In this thesis, \textit{Urgency Index} is therefore treated as a simple proxy for the pressure created by the reveal message.

The third dimension is execution clarity. A signal may attract attention and create pressure, but participants still need to know what to do. Clear trading instructions can reduce practical friction by naming the target, exchange, pair, or expected action. This does not mean that longer instructions are always better. A very long message may slow interpretation, while a very short message may leave practical uncertainty. The key theoretical point is that instruction content can affect how easily attention becomes actual order flow. In the thesis, \textit{Instruction Density} captures this practical side of the signal.

These three variables are proxies, not direct measures of trader psychology. They do not observe private beliefs, individual order submission, or the exact moment at which each member reads the message. Their role is more limited and more defensible. They translate observable message features into coordination mechanisms supported by the literature. \textit{Hype Intensity} relates to attention before the reveal, \textit{Urgency Index} relates to compressed decision time at the reveal, and \textit{Instruction Density} relates to the ease of turning the signal into action. This mapping is important because it connects the empirical treatment variables to theory rather than treating them as arbitrary text features.

These dimensions do not assume that every participant believes the signal or that the asset has changed in value. They instead treat message style as a way to coordinate mass behavior. Countdown language, urgent tone, and explicit instructions can focus attention, reduce hesitation, and make the buying race easier to join. Research on openly declared cryptocurrency pumps also shows that some traders know the risks but still gamble on exiting before the price falls \cite{dhawan2023new}. In this setting, the signal becomes a race trigger rather than a claim about long-term value, which supports studying the relationship between message characteristics and short-term price behavior.

\section{Market Microstructure, Causal Measurement, and Outcome Window}
\label{sec:market-microstructure-short-horizon}
\label{sec:detection-to-causal-measurement}

Pump-and-dump schemes do not operate in an abstract market. Their effect depends on the market structure of the asset being targeted. Market microstructure research studies how trading mechanisms, order flow, liquidity, spreads, and information affect prices over short horizons. This perspective is useful because a pump signal works through immediate trading pressure. The question is not only whether people receive the message, but also whether the market can absorb the orders that follow.

Classic market microstructure theory shows that order flow can move prices when trades carry information or when liquidity is limited \cite{kyle1985continuous}. Even when a trade does not contain true fundamental information, the market may still react to the buying pressure. In a thin market, the next available sell orders may be placed at much higher prices, so a small amount of demand can create a large return. This mechanism is directly relevant to cryptocurrency pumps, where many target assets have limited depth and uneven trading activity.

Illiquidity is therefore a key concept. The Amihud illiquidity measure captures how strongly prices move in response to trading volume \cite{amihud2002illiquidity}. Cryptocurrency evidence also shows that higher liquidity is linked to weaker return predictability and stronger market efficiency \cite{wei2018liquidity}. In the context of pump-and-dump schemes, this idea means that the same coordinated buying wave can have a much larger impact in a thin asset than in a liquid one. A large asset with deep order books can absorb more demand. A small or thinly traded asset may react sharply to much less volume.

High-frequency market measures also matter because pump events happen in minutes. Short-run volatility, bid-ask spread, zero-return intervals, and abnormal volume can all describe the trading environment before a signal. The Roll spread estimator provides a classic way to infer effective bid-ask spreads from price changes \cite{roll1984simple}. Range-based volatility estimators also show that high and low prices within a short interval contain information beyond close-to-close returns \cite{parkinson1980extreme}. Pump-and-dump studies further show that abnormal price and volume movements are central signs of these events \cite{kamps2018pump, li2021cryptocurrency}. These measures support the broader idea that short-horizon market conditions are not merely noise.

This is especially relevant for event studies in crypto markets. A large return immediately after a pump signal may come from the signal itself, but it may also reflect a market that was already unstable, illiquid, or moving prior to the reveal. Abnormal volume before $T_0$ may indicate early activity, while high pre-event volatility may show that the coin was already risky. A careful empirical design therefore needs to observe the market before the signal, not only after it.

Cryptocurrency markets make this issue stronger. Trading is global and continuous, and the same asset may trade across several exchanges at different prices \cite{makarov2020trading}. Some tokens have active markets, while others have long periods with little or no trading. In the thesis sample, many coins also trade against BTC (Bitcoin), ETH (Ether), BNB (Binance Coin), or different cryptocurrency coins, which creates additional variation in how returns are measured and understood. These conditions make it especially important to separate the effect of the pump signal from the market state that existed before the signal.

The literature on cryptocurrency pump-and-dump schemes confirms that price, volume, and volatility can rise sharply after an announcement, followed by rapid reversals \cite{li2021cryptocurrency}. Studies that observe the full pump ecosystem also show that less popular or less liquid coins are more sensitive to these events \cite{hamrick2021examination}. This evidence fits the market microstructure logic. Manipulation is easier when the market is shallow enough for coordinated demand to move the price.

Market capitalization is not identical to liquidity, but it is a useful way to connect this theory to empirical analysis. Larger assets tend to have broader attention, deeper trading activity, and more resistance to sudden demand shocks. Smaller assets are more exposed to sharp moves. In this thesis, market capitalization is used as the main practical vulnerability proxy because it captures an important part of this difference and is available consistently across the event sample \cite{amihud2002illiquidity, hamrick2021examination}.

Target selection is therefore not random. Organizers have incentives to choose coins that combine several conditions. The asset must be known enough for members to recognize it quickly, listed on an exchange where members can trade, and active enough for orders to be executed. At the same time, it must not be so liquid that the buying wave disappears into a deep order book. Recent evidence on pump-and-dump drivers also supports the view that both coin-level and market-level factors matter, including market capitalization, trading volume, social-media attention, and exchange availability \cite{charfeddine2024drives}. This fits the idea that pump targets are selected from a narrow part of the crypto market, not from all assets equally.

This creates an important boundary condition for the thesis. Lower market capitalization can make a coin easier to move \cite{amihud2002illiquidity, hamrick2021examination}, but the smallest assets may also trade so rarely that the measured reaction becomes unstable. A very inactive coin may show stale prices, zero-volume intervals, or large jumps caused by a small number of trades. Such conditions are part of the market vulnerability, but they also make causal measurement harder. The ideal target is therefore not simply the smallest coin. It is a coin that is thin enough for coordinated demand to matter, but active enough for the event to appear in market data.

Pre-event market conditions also matter because they can influence both the organizer's choices and the later price response \cite{charfeddine2024drives, hu2023sequence}. A coin with unusual volume before the reveal may already be attracting attention or may reflect early activity by better-informed traders \cite{li2021cryptocurrency}. A coin with high volatility may be more likely to move sharply even without the signal \cite{parkinson1980extreme}. A coin with wide effective spreads \cite{roll1984simple} or frequent zero-return intervals may react differently from a coin with smoother trading. These conditions help explain why the same signal style can produce different outcomes across assets.

This logic also shows why market capitalization should be interpreted carefully. It is not a direct order-book measure and cannot fully describe depth \cite{amihud2002illiquidity}, spread \cite{roll1984simple}, or active trader participation. Two coins with similar capitalization may differ in exchange structure, quoted pair \cite{makarov2020trading}, trading activity \cite{wei2018liquidity}, and recent volatility. Still, market capitalization remains useful as a broad proxy for the asset's resistance to coordinated demand \cite{charfeddine2024drives}. It supports the theoretical expectation that signal effects are conditional on the market environment rather than uniform across all target coins.

This microstructure perspective also explains why the thesis focuses on very short horizons. Over a long period, many other factors can affect returns, including market-wide news, changes in Bitcoin or Ethereum prices, exchange-specific events, or normal investor activity \cite{mackinlay1997event}. In the first minutes after a reveal, the direct connection between the message and the market reaction is much tighter \cite{xu2019anatomy, li2021cryptocurrency}. The short window therefore matches both the behavior of pump groups and the trading mechanism through which their signals work.

This market setting also shapes how causal measurement should be defined. Existing research has largely focused on detection and target prediction \cite{kamps2018pump, victor2019cryptocurrency, lamorgia2020pump, hu2023sequence}. Detection asks whether a pump-and-dump event can be found in market or message data. It usually relies on abnormal price and volume movements, suspicious timing, or social-media activity. This work is useful because it shows that manipulation leaves measurable traces, but it does not by itself explain which part of the signal created the price response.

Target-coin prediction asks whether the coin in a scheduled pump can be predicted before the public announcement \cite{hu2023sequence}. This line of research is valuable because it studies information leakage, pre-pump behavior, and organizer choices. It also shows that pump events are not random shocks to the market. A model may learn organizer habits, channel histories, or pre-event anomalies, but those patterns do not show why the final public signal generated a specific return.

A third task, outcome prediction, asks whether an event is likely to produce a positive standardized short-window return. This definition is useful for comparing events, but it is not the same as a realistic trading profit \cite{mackinlay1997event}. A trader would also face entry delay, exit timing, fees, spreads \cite{amihud2002illiquidity}, and slippage \cite{kyle1985continuous}. The prediction task therefore classifies a standardized event outcome, not whether a participant could safely earn money from the scheme.

Causal measurement asks a different question. It asks whether signal characteristics, such as urgency, explain post-reveal returns after relevant background conditions are considered \cite{rubin1974estimating, angrist2009mostly}. This distinction is important because a variable can help classify events without being the mechanism that caused the market reaction \cite{shmueli2010explain}. For example, a channel may repeatedly choose certain coins, or organizers may use stronger wording when they already expect a strong response. Such variables can be useful for prediction while still being difficult to interpret causally.

Isolating signal effects is challenging because pump signals are not issued at random. Organizers often select coins that are already volatile, thinly traded \cite{charfeddine2024drives}, or showing unusual activity before the event \cite{li2021cryptocurrency}. They may also craft more aggressive messages when they expect a coin to move strongly. Strong wording can therefore be part of the reaction mechanism and, at the same time, a sign of organizer expectations. Without proper adjustment, it is easy to mistake a coin's pre-existing market state for the direct effect of a signal \cite{angrist2009mostly}.

To address this, the study adopts an event-level perspective, focusing on the pump session rather than isolated messages. This approach aligns with the real-world organization of these schemes. Preparatory communication builds attention, the target coin is revealed, and the market reacts.

The causal framework relies on a simple conceptual distinction. Signal characteristics are the candidate drivers of interest, the short-horizon price response is the outcome to be explained, and pre-existing market or organizer conditions are possible confounders.

This approach follows the potential outcomes view of causality \cite{rubin1974estimating}. In this context, we can imagine two possible paths for a single event, namely the observed return following the actual signal and the return that would have occurred had the signal been delivered in a different style. Since only one path is ever observable, the empirical task is to construct a credible comparison through observational data \cite{angrist2009mostly}.

This comparison is credible only to the extent that relevant pre-reveal differences are observed and adjusted for. The main challenge is that organizers choose both the target coin and the signal style. Stronger language may influence participants, but it may also reflect the organizer's private expectations about the event. The empirical design therefore has to compare events that are similar in observed market and organizer-related conditions while differing in signal characteristics. This does not remove every possible source of bias, but it makes the causal question more precise.

The outcome window has to match how these schemes actually work. Past studies show that pump reactions happen extremely fast. Research on the anatomy of cryptocurrency pump-and-dump schemes reports that the price usually reaches its peak just a few minutes after the coin announcement \cite{xu2019anatomy}. Broader pump-and-dump samples also show sharp short-term increases in price, volume, and volatility, followed by quick reversals \cite{li2021cryptocurrency}.

This is why the thesis focuses on a 10-minute window. One major empirical study analyzes market activity in 10-second intervals up to 600 seconds after a pump announcement and shows that even after ten minutes, price, trading volume, and volatility remain significantly higher than before the announcement \cite{li2021cryptocurrency}. Based on this timing evidence, this thesis measures return at the 10-minute mark as a standardized checkpoint within the immediate post-reveal period.

A natural alternative might be to look for the maximum price reached during a broader window, such as the first 20 minutes. A fixed 10-minute mark provides a stronger basis for comparison because it measures the same phase of the manipulation in every event. It is also more reliable than using a maximum price, which may reflect only a brief trade in a thin market rather than a broader market reaction. A fixed checkpoint better reflects trading reality, since traders do not know in advance when the peak will occur.

This outcome choice also follows the broader logic of event studies. Event-study methods measure the market response close to a defined event time, since the effect of the event should be most visible near the announcement when prices adjust quickly \cite{mackinlay1997event}. In the pump setting, the event time is unusually clear because the public reveal defines $T_0$. This makes a short fixed window more appropriate than a long horizon in which unrelated market news and normal trading conditions can enter the return.

Another possible approach would be to compare the target coin with a synthetic counterfactual and measure cumulative abnormal returns. Synthetic-control methods are useful when a weighted control group can closely track the treated unit before the event \cite{abadie2010synthetic, abadie2021synthetic}. In sparse cryptocurrency markets, however, many pump targets do not move closely with broad market benchmarks. If the counterfactual path is weak, subtracting it can add noise instead of removing it.

The 10-minute window is also short enough to limit the influence of unrelated market factors that become more important over longer horizons \cite{li2021cryptocurrency, lamorgia2023doge}. Looking at hourly or daily returns would be less appropriate, because pump-and-dumps do not reflect changes in fundamental value. Once the scheme ends, the targeted coin typically experiences a sharp price drop as the artificial inflation collapses \cite{kamps2018pump}. Staying within the first 10 minutes therefore reduces exposure to later unrelated market movements and keeps the analysis focused on the immediate response to the coordinated signal.

Even though the 10-minute window reduces later market noise, it does not fully remove the influence of what was happening to the coin before the pump began. That is why the study still needs to control for the market conditions just before the signal was posted. The aim is not to eliminate all randomness, but to compare events that started from similar conditions.

These identification issues explain why the empirical framework combines causal estimation with predictive modeling. The causal component asks whether signal characteristics are associated with short-horizon returns after accounting for relevant background conditions \cite{angrist2009mostly, hu2023sequence}. The predictive component asks whether information available up to the reveal can classify events by their likely short-term outcome. As \cite{shmueli2010explain} argues, explanation and prediction require different modeling choices and different interpretations, so the two tasks are kept separate.

This framework also connects the literature to the data design. The literature shows that pumps are fast, socially coordinated, and shaped by liquidity \cite{xu2019anatomy, li2021cryptocurrency, hamrick2021examination}. The data therefore has to capture the message, the reveal time, the pre-reveal market state, and the immediate post-reveal return. The methodology then turns this structure into two empirical questions. These questions ask what explains the reaction and what predicts the short-term outcome.

\section{Theoretical Expectations and Hypotheses}
\label{sec:theoretical-expectations-hypotheses}

The hypotheses for this study are derived from three interconnected theoretical premises. They concern how the message influences trader behavior, how market conditions constrain price movements, and how these factors combine to make the short-term outcome predictable. The literature reviewed above shows that pump-and-dump events are not random price spikes only. They are organized communication events that take place in markets with different levels of depth and liquidity.

The first hypothesis focuses on how coordinated messages facilitate group action through attention-building, urgency, and practical trading instructions. Since organizers aim to trigger immediate buying, they may build anticipation before the reveal, use forceful wording at the reveal, and provide instructions that reduce confusion among participants \cite{lamorgia2023doge, xu2019anatomy}. These signal characteristics matter because a pump is a speed-based event. A trader who reacts slowly may enter after the price has already moved, while a trader who reacts early may benefit from the first wave of buying pressure.

The retained signal characteristics can work through several channels. \textit{Hype Intensity} can concentrate attention before the target is known. \textit{Urgency Index} can make delay feel costly and reduce the space for careful evaluation. \textit{Instruction Density} can reflect how much practical trading guidance is contained in the signal. Together, these features can create a shared belief that other participants are about to buy immediately. This kind of belief is powerful in a coordination setting. Traders may buy not because they believe the asset is valuable, but because they believe many other traders will buy at the same moment.

The first hypothesis states that stronger values of the retained signal-characteristic vector, namely \textit{Hype Intensity}, \textit{Urgency Index}, and \textit{Instruction Density}, are expected to have positive estimated effects on short-term cryptocurrency returns after the coin reveal.

These signal characteristics may not be effective in every market environment. If they have measurable effects, those effects are expected to appear mainly in the immediate post-reveal period, before the initial momentum fades or sellers begin to exit.

The second hypothesis focuses on how the targeted asset changes the effect of the signal. The same signal is unlikely to produce identical results across different coins. A large, liquid market can absorb strong buying pressure with limited price movement, while a thin market is much more sensitive. This logic is motivated by literature on market manipulation and illiquidity \cite{amihud2002illiquidity, dhawan2023new} and by evidence that lower-popularity coins can react more strongly in pump-and-dump events \cite{hamrick2021examination}.

Market capitalization is central to this expectation. It does not perfectly measure order-book depth, but it captures an important part of the difference between widely traded assets and smaller, less active ones. A small-cap coin can be visible enough for traders to recognize it, while still being thin enough for coordinated demand to move the price.

The second hypothesis states that the estimated effects of \textit{Hype Intensity}, \textit{Urgency Index}, and \textit{Instruction Density} are expected to be stronger in lower-cap assets than in higher-cap assets, and to weaken as market capitalization increases.

This does not imply that all low-cap assets are easily manipulated. Rather, it suggests that market size may change how strong the signal effect is. Sparse trading can make measurement difficult at the lowest end of the market, while deeper markets may reduce the effect of the same signal at the higher end.

The final hypothesis addresses predictability. If the short-term outcome of a pump depends on signal characteristics and pre-event market states, then those variables should contain useful information before the outcome is realized. Prior detection and prediction studies show that pump events leave observable traces in market and communication data \cite{kamps2018pump, lamorgia2020pump, hu2023sequence}. This supports the expectation that the same information environment can help classify short-term outcomes.

The third hypothesis states that observable market patterns up to the reveal, combined with the same signal characteristics, are expected to contain information that can help predict short-term pump outcomes.

While the first two hypotheses focus on isolating causal drivers, the third hypothesis is a forecasting claim and must be interpreted separately from the causal results.

Having established the theoretical motivation and hypotheses, the next chapter turns to the empirical foundation of the study: the identification of pump events, their connection to market data, and the construction of the variables used in the analysis.
